\subsection{An alternative coordinate character calculation}
\label{sec:coordinate-character-alternative}
The following direct calculation is independent of the character decomposition in Lemma~\ref{lem:orthogonal-exterior-cube}; it is not needed for the main proof.
Use an orthonormal basis and write \(\rho=\sum_i x_i^2\).

Any possible image of \(\rho^2\), with its
determinant twist, must transform by the determinant character of
\(O_4\) in \(\bigwedge^3\Sym^2V\).  There is no such vector.
To see this directly, invariance with this character under the four
coordinate reflections requires each variable to have odd total
degree in any nonzero exterior monomial.  Its total degree is six,
so the multidegree must be a permutation of \((3,1,1,1)\).
For a fixed variable \(x_i\) of degree three, the possible
monomials are
\[
 \omega_i=\bigwedge_{j\ne i}x_ix_j,
 \qquad x_i^2\wedge(x_ix_j)\wedge(x_kx_l)
       \quad(\{i,j,k,l\}=\{1,2,3,4\}).
\]
The coefficient of the second kind is zero: the transposition of
\(x_k,x_l\) fixes that monomial but has determinant minus one.
Only a combination of the \(\omega_i\) remains.  Applying the
infinitesimal rotation
\(x_i\partial/\partial x_j-x_j\partial/\partial x_i\),
the coefficient of
\(x_i^2\wedge(x_ix_k)\wedge(x_ix_l)\) comes, up to sign,
only from \(\omega_i\).  Invariance forces its coefficient to
vanish, for every \(i\).  This proves the assertion.

